\documentclass[letterpaper]{article}

\usepackage{aaai2027}
\usepackage{times}
\usepackage{helvet}
\usepackage{courier}
\usepackage{algorithm}
\usepackage{algorithmic}
\usepackage{booktabs}
\usepackage{amsmath}

\title{Supplementary Material }

\author{}

\begin{document}

\maketitle


\section{Additional Method Details}
\label{sec:supp-method}

This supplement provides additional implementation details and the complete
inference procedure of the latent memory manager.


\subsection{Implementation Details}
\label{sec:supp-implementation}

The latent memory manager is implemented as an inference-only extension
to the original MemGen pipeline. The Trigger, Weaver, and Reasoner
parameters remain unchanged throughout evaluation, and no additional
training, parameter updates, or optimization objectives are introduced.

Each memory slot stores a detached Weaver-space latent memory, its
retrieval key, and the associated last-retrieval counter. Before insertion,
each generated latent memory is detached from the computation graph and
transferred to CPU storage. The retrieval key is computed from the stored
latent representation after detachment:

\[
k_t =
\operatorname{MeanPool}(\widetilde{m}_t).
\]

The memory bank is initialized independently for each context and destroyed
after the corresponding inference session. It does not share memory states
across contexts. Each stored memory consists of eight latent-token vectors
generated by the Weaver. The query representation uses mean pooling over
the most recent \(L=64\) Reasoner hidden states.

All experiments use batch size one and bfloat16 inference. Decoding is
performed greedily with a fixed generation budget of 40 tokens.


\subsection{Model Configuration}
\label{sec:supp-model}

% TODO:
% Base model checkpoint
% MemGen version
% Trigger / Weaver / Reasoner architecture
% Hidden dimension
% Latent token configuration
% Model loading configuration


\subsection{Hardware and Runtime Configuration}
\label{sec:supp-runtime}

% TODO:
% GPU model
% CUDA version
% PyTorch version
% inference environment


\subsection{Latent Memory Bank Configuration}
\label{sec:supp-bank-config}

Unless otherwise stated, the latent memory manager uses the following
default configuration:

\begin{table}[h]
\centering
\begin{tabular}{lc}
\toprule
Component & Configuration\\
\midrule
Memory capacity $C$ & 16 slots\\
Retrieval top-$k$ & 2\\
Retrieval threshold $\tau_r$ & 0.05\\
Update threshold $\tau_u$ & 0.10\\
Temporal decay $\alpha$ & 0.05\\
Query pooling length $L$ & 64\\
Stored latent tokens per slot & 8\\
Storage device & CPU\\
Precision & bfloat16\\
\bottomrule
\end{tabular}
\caption{Default configuration of the latent memory manager.}
\label{tab:supp-bank-config}
\end{table}



\subsection{Complete Inference Procedure}
\label{sec:supp-algorithm}

The complete inference procedure, including retrieval, latent-memory
generation, and memory update operations, is summarized in
Algorithm~\ref{alg:latent-memory}.


\begin{algorithm}[t]
\caption{Inference procedure of the latent memory manager}
\label{alg:latent-memory}
\begin{algorithmic}[1]

\STATE \textbf{Input:} hidden states \(H_1,\ldots,H_T\), retrieval
threshold \(\tau_r\), update threshold \(\tau_u\), temporal decay
coefficient \(\alpha\), maximum retrieval number \(K\), memory capacity
\(C\)

\STATE Initialize memory bank
\(\mathcal{M}\leftarrow\varnothing\)
and retrieval counter \(r\leftarrow0\)

\FOR{each reasoning step \(t=1,\ldots,T\)}

    \STATE \(g_t\leftarrow\operatorname{Trigger}(H_t)\)

    \IF{\(g_t=\texttt{INVOKE}\)}

        \STATE
        \(q_t\leftarrow\operatorname{MeanPool}(H_t[-L:])\)

        \STATE
        \(r\leftarrow r+1\)

        \FOR{each memory \(m_i\in\mathcal{M}\)}

            \STATE Compute retrieval score:
            \[
            s_i=
            \operatorname{cosine}(q_t,k_i)
            \exp(-\alpha\Delta r_i)
            \]

        \ENDFOR


        \STATE Retrieve:
        \[
        R_t=
        \operatorname{TopK}
        (\{m_i:s_i\geq\tau_r\},K)
        \]


        \FOR{each \(m_i\in R_t\)}

            \STATE Update last retrieval count:
            \(r_i^{\mathrm{last}}\leftarrow r\)

        \ENDFOR


        \STATE Generate latent memory:
        \[
        m_t\leftarrow
        \operatorname{Weaver}(H_t,R_t)
        \]


        \STATE Update Reasoner:
        \[
        y_t\leftarrow
        \operatorname{Reasoner}(H_t,m_t)
        \]


        \STATE Detach and transfer:
        \[
        \widetilde{m}_t
        \leftarrow
        \operatorname{detach}(m_t).\operatorname{cpu}()
        \]


        \STATE Compute retrieval key:
        \[
        k_t
        \leftarrow
        \operatorname{MeanPool}(\widetilde{m}_t)
        \]


        \IF{\(\mathcal{M}=\varnothing\)}

            \STATE Insert
            \((\widetilde{m}_t,k_t,r)\)
            into \(\mathcal{M}\)


        \ELSE

            \STATE Find closest existing slot:
            \[
            j\leftarrow
            \arg\max_i
            \operatorname{cosine}(k_t,k_i)
            \]


            \STATE Compute update similarity:
            \[
            u_t=
            \operatorname{cosine}(k_t,k_j)
            \]


            \IF{\(u_t\geq\tau_u\)}

                \STATE Replace slot \(j\) with
                \((\widetilde{m}_t,k_t,r)\)


            \ELSE

                \IF{\(|\mathcal{M}|=C\)}

                    \STATE Select eviction slot:
                    \[
                    e\leftarrow
                    \arg\max_i
                    (r-r_i^{\mathrm{last}})
                    \]

                    \STATE Evict slot \(e\)


                \ENDIF


                \STATE Insert
                \((\widetilde{m}_t,k_t,r)\)
                into \(\mathcal{M}\)


            \ENDIF

        \ENDIF


    \ELSE

        \STATE
        \(y_t\leftarrow\operatorname{Reasoner}(H_t)\)


    \ENDIF

\ENDFOR

\end{algorithmic}
\end{algorithm}

\section{Additional Baseline Details}
\label{sec:supp-baselines}

This section provides additional implementation details of the explicit
textual-memory baselines used in the experiments.


\subsection{Dense Retrieval Baseline}
\label{sec:supp-dense}

The dense retrieval baseline uses an external embedding-based retrieval
pipeline. Contexts are divided into the same text chunks used by the
corresponding textual-memory comparison. Each chunk is encoded into a dense
vector representation, and retrieval is performed using cosine similarity
between the query representation and stored chunk embeddings.

% TODO:
% Specify embedding model name and version.
% Specify chunk size and indexing details.


\subsection{Textual Memory Controls}
\label{sec:supp-text-memory}

The textual-memory controls evaluate whether explicit text retrieval can
provide comparable benefits under matched evaluation conditions.

The rolling-summary baseline maintains a compressed textual summary of the
processed context and provides this summary during question answering.

The BM25 retrieval baseline constructs a sparse retrieval index over context
chunks and retrieves the top-ranked chunks for each question.

The matched-budget textual control restricts retrieved textual information
to a fixed token budget comparable to the latent-memory representation.

% TODO:
% Add exact summary length, BM25 implementation, and indexing details.
\section{Additional Ablation and Sensitivity Studies}
\label{sec:supp-ablation}

This section reports additional ablations and sensitivity analyses
for the latent memory manager on EventQA.


\subsection{Shared Evaluation Protocol}
\label{sec:supp-protocol}

Unless explicitly varied, all experiments use the default latent memory
manager configuration.

The default setting uses 4,096-token construction chunks, a memory capacity
of \(C=16\), query retrieval top-\(k=2\), retrieval threshold \(0.05\),
construction update threshold \(0.10\), temporal decay coefficient
\(\alpha=0.05\), and a 40-token generation budget.

Query-time writes are disabled, and the constructed memory bank snapshot is
restored before each question. Each condition is evaluated over five
complete process-level passes over five contexts and 500 questions, using
base seeds 42, 142, 242, 342, and 442.

Results are reported as mean \(\pm\) population standard deviation.


\subsection{Capacity Sensitivity}
\label{sec:supp-capacity}

We vary the maximum memory bank capacity while keeping all other components
unchanged.

\begin{table}[t]
\centering
\begin{tabular}{lccc}
\toprule
Capacity & EM & Recall & Format failures\\
\midrule
$C=8$ &
0.156$\pm$0.035 &
0.209$\pm$0.027 &
136.4$\pm$13.0\\

$C=16$ &
\textbf{0.188$\pm$0.047} &
\textbf{0.231$\pm$0.042} &
\textbf{118.6$\pm$19.1}\\

$C=24$ &
0.190$\pm$0.027 &
0.228$\pm$0.023 &
125.8$\pm$24.7\\
\bottomrule
\end{tabular}
\caption{Bank-capacity sensitivity on EventQA.}
\label{tab:supp-capacity}
\end{table}


Reducing capacity to \(C=8\) lowers EM and recall while increasing format
failures. This configuration reaches capacity in every context. The
\(C=24\) setting constructs only 16--17 slots in practice and therefore
does not demonstrate a benefit beyond the realized occupancy of the default
configuration.


\subsection{Query Retrieval Depth Sensitivity}
\label{sec:supp-topk}

We evaluate different numbers of retrieved latent memories during query-time
retrieval.

\begin{table}[t]
\centering
\begin{tabular}{lccc}
\toprule
Top-$k$ & EM & Recall & Format failures\\
\midrule
$k=1$ &
0.050$\pm$0.006 &
0.220$\pm$0.013 &
338.0$\pm$31.0\\

$k=2$ &
\textbf{0.188$\pm$0.047} &
\textbf{0.231$\pm$0.042} &
\textbf{118.6$\pm$19.1}\\

$k=4$ &
0.128$\pm$0.029 &
0.143$\pm$0.033 &
181.0$\pm$67.8\\
\bottomrule
\end{tabular}
\caption{Query retrieval depth sensitivity.}
\label{tab:supp-topk}
\end{table}


Retrieving only one memory substantially reduces EM and increases format
failures. Increasing retrieval depth to four memories also underperforms the
default \(k=2\) configuration, suggesting that a small number of relevant
latent memories provides a better balance.


\subsection{Query Retrieval Threshold Sensitivity}
\label{sec:supp-threshold}

We vary the threshold controlling whether stored memories participate in
query-time retrieval.

\begin{table}[t]
\centering
\begin{tabular}{lccc}
\toprule
Threshold & EM & Recall & Format failures\\
\midrule
0.03 &
0.176$\pm$0.023 &
\textbf{0.242$\pm$0.009} &
155.2$\pm$27.2\\

0.05 &
\textbf{0.188$\pm$0.047} &
0.231$\pm$0.042 &
\textbf{118.6$\pm$19.1}\\

0.10 &
0.008$\pm$0.000 &
0.178$\pm$0.000 &
377.0$\pm$0.0\\
\bottomrule
\end{tabular}
\caption{Query retrieval-threshold sensitivity.}
\label{tab:supp-threshold}
\end{table}


A lower threshold slightly improves recall but reduces EM and format
stability. At threshold \(0.10\), no candidate memory passes the retrieval
criterion across the evaluated questions, effectively removing retrieved
latent-memory conditioning.

\subsection{Construction Update Threshold Sensitivity}
\label{sec:supp-update-threshold}

The construction update threshold determines whether a newly generated latent
memory refreshes an existing thread or creates a new thread. Query-time
retrieval remains fixed at threshold \(0.05\) and top-\(k=2\).

\begin{table}[t]
\centering
\begin{tabular}{lccc}
\toprule
Update threshold & EM & Recall & Format failures\\
\midrule
0.00 &
0.153$\pm$0.037 &
0.215$\pm$0.019 &
192.2$\pm$36.5\\

0.05 &
0.147$\pm$0.060 &
0.193$\pm$0.051 &
164.4$\pm$67.3\\

0.10 &
\textbf{0.188$\pm$0.047} &
\textbf{0.231$\pm$0.042} &
\textbf{118.6$\pm$19.1}\\
\bottomrule
\end{tabular}
\caption{Construction update-threshold sensitivity.}
\label{tab:supp-update-threshold}
\end{table}


Lower update thresholds cause more generated memories to overwrite existing
threads and underperform the default configuration. Within the tested range,
the \(0.10\) setting provides the strongest overall balance.


\subsection{New-Thread Append Ablation}
\label{sec:supp-append}

This ablation suppresses unmatched new-thread insertion during construction
while preserving retrieval, matched-thread replacement, capacity management,
and frozen query evaluation.

\begin{table}[t]
\centering
\begin{tabular}{lccc}
\toprule
Construction policy & EM & Recall & Format failures\\
\midrule
No new-thread append &
0.052$\pm$0.009 &
0.215$\pm$0.010 &
350.2$\pm$14.0\\

Default configuration &
\textbf{0.188$\pm$0.047} &
\textbf{0.231$\pm$0.042} &
\textbf{118.6$\pm$19.1}\\
\bottomrule
\end{tabular}
\caption{New-thread append ablation.}
\label{tab:supp-append}
\end{table}


Suppressing new-thread creation discards historical latent information that
cannot be matched to existing slots. The resulting degradation indicates
that retaining distinct latent-memory threads is important under the
evaluated protocol.


\subsection{Frozen-Bank Protocol Diagnostic}
\label{sec:supp-frozen}

The default evaluation freezes the constructed context bank and restores the
same snapshot before every question. We additionally evaluate a mutable
sequential-query variant that retains one bank throughout the ordered
questions of each context and enables query-time updates.

\begin{table}[t]
\centering
\begin{tabular}{lccc}
\toprule
Protocol & EM & Recall & Format failures\\
\midrule
Mutable sequential queries &
0.140$\pm$0.038 &
0.196$\pm$0.036 &
167.4$\pm$24.8\\

Frozen default &
\textbf{0.188$\pm$0.047} &
\textbf{0.231$\pm$0.042} &
\textbf{118.6$\pm$19.1}\\
\bottomrule
\end{tabular}
\caption{Frozen-bank protocol diagnostic.}
\label{tab:supp-frozen}
\end{table}


The mutable variant introduces question-derived writes during evaluation
and produces lower EM and recall with more format failures. This result is
consistent with query-induced state drift and supports the frozen-bank
evaluation protocol used in the main experiments.



\section{Additional Experimental Details}
\label{sec:supp-analysis}


\subsection{Additional Efficiency Measurements}

% TODO:
% Detailed runtime breakdown.
% GPU memory allocation analysis.
% Additional efficiency measurements.


\subsection{Additional Error Analysis}

% TODO:
% Successful retrieval examples.
% Failure cases.
% Parser-format failure analysis.


\end{document}